\documentclass[../../../main.tex]{subfiles}
\begin{document}
\subsection{Creation and Annihilation Operator}
The operators $\hat{a}_{\mathbf{k }s}$ and $\hat{a}^\dagger_{\mathbf{k }s}$ are called annihilation and creation operators because they respectively destroy and create one photon with wavevector $\mathbf{k}$ and polarization $s$, namely
\begin{align*}
  \hat{a}_{\mathbf{k }s} | \cdots n_{\mathbf{k }s} \cdots \rangle         & = \sqrt{n_{\mathbf{k }s}} |  \cdots n_{\mathbf{k }s} - 1 \cdots \rangle    \\
  \hat{a}^\dagger_{\mathbf{k }s} | \cdots n_{\mathbf{k }s} \cdots \rangle & = \sqrt{n_{\mathbf{k }s} + 1} |\cdots  n_{\mathbf{k }s} + 1 \cdots \rangle
\end{align*}
Consequently, for the vacuum state $|0\rangle$ one define
\begin{align*}
  \hat{a}_{\mathbf{k }s} | 0 \rangle   & = 0_F                                                 \\
  \hat{a}^+_{\mathbf{k }s} | 0 \rangle & = | 1_{\mathbf{k }s} \rangle = | \mathbf{k} s \rangle
\end{align*}
where $0_F$ is the zero of the Fock space (usually indicated with $0$), and $| k s \rangle$ is clearly the state of one photon with wavevector $k$ and polarization $s$, such that
\begin{equation*}
  \langle \mathbf{r} | \mathbf{k} s \rangle = \frac{e^{i \mathbf{k} \cdot \mathbf{r}}}{\sqrt{V} \epsilon_{\mathbf{k }s}}
\end{equation*}
Any higher Foch state then can be written as 
\begin{equation*}
  \ket{n_{\mathbf{k }s}}=\frac{1}{\sqrt{n!}}\left( \hat{a}^\dagger_{\mathbf{k }s}   \right) ^{n }\ket{0}
\end{equation*}

It follows immediately that the number operator
\begin{equation*}
  \hat{N }_{\mathbf{k }s } = \hat{a } ^\dagger_{\mathbf{k }s }\hat{a }_{\mathbf{k }s}
\end{equation*}
have the action of
\begin{equation*}
  \hat{N}_{\mathbf{k }s} | \cdots n_{\mathbf{k }s} \cdots \rangle = n_{\mathbf{k }s}  | \cdots n_{\mathbf{k }s}  \cdots \rangle
\end{equation*}
The follow-up is as follows
\begin{align*}
  \hat{a}^\dagger_{\mathbf{k }s} \hat{a}_{\mathbf{k }s} | \cdots n_{\mathbf{k}s} \cdots \rangle
   & =  \hat{a }^\dagger_{\mathbf{k }s}\sqrt{n_{\mathbf{k}s}} | \cdots (n_{\mathbf{k}s} - 1) \cdots \rangle \\
  \hat{N }_{\mathbf{k }s}| \cdots n_{\mathbf{k}s} \cdots \rangle
   & =  n_{\mathbf{k}s} | \cdots n_{\mathbf{k}s} \cdots \rangle
\end{align*}

We can also prove the relation
\begin{equation*}
  \sqrt{n! }  \ket{0 }= \left( \hat{a}^\dagger_{\mathbf{k }s}   \right) ^n \ket{0}
\end{equation*}
using induction. 
For the base case, we use $n=0$, which according to the hypothesis
\begin{equation*}
  \left( \hat{a}^\dagger_{\mathbf{k }s}  \right) ^0 \ket{0 }=0! \ket{0}
\end{equation*}
is true.
Then for the $n+1$
\begin{equation*}
  \left( \hat{a}^\dagger_{\mathbf{k }s}   \right) ^{n+1 }\ket{0 }=\hat{a}^\dagger_{\mathbf{k }s}  \left( \hat{a}^\dagger_{\mathbf{k }s}   \right) ^n \ket{0 }
\end{equation*}
On using the hypothesis
\begin{equation*}
  \hat{a}^\dagger_{\mathbf{k }s} \sqrt{n_{\mathbf{k }s}! } \ket{n_{\mathbf{k }s} }\sqrt{n_{\mathbf{k }s}! }\sqrt{n_{\mathbf{k }s}+1}\ket{n_{\mathbf{k }s}+1}= \sqrt{(n_{\mathbf{k }s}+1)!} \ket{n_{\mathbf{k }s}+1}
\end{equation*}
we find that it holds true.

\subsubsection{Derivation: the creation and annihilation operator}
With $\hat{N}= \hat{a}^\dagger_{\mathbf{k }s}  \hat{a}_{\mathbf{k }s} $ as the number operator, we want to observe what is the action of $\hat{a}_{\mathbf{k }s} $ and $\hat{a}^\dagger_{\mathbf{k }s} $ to the eigenstate of $\hat{N}$.
That is we consider the state $\hat{a}_{\mathbf{k }s}  \ket{n_{\mathbf{k }s}}$ using said operator
\begin{align*}
    \hat{N} \left( \hat{a}_{\mathbf{k }s} \ket{n_{\mathbf{k }s}} \right) &=  \hat{a}_{\mathbf{k }s} (\hat{N }-1  ) \ket{n_{\mathbf{k }s}}  \\
    \hat{N} \left( \hat{a}_{\mathbf{k }s} \ket{n_{\mathbf{k }s}} \right) &=  (n_{\mathbf{k }s}-1  ) \left( \hat{a}_{\mathbf{k }s} \ket{n_{\mathbf{k }s}} \right) 
\end{align*}
As such, the operator $\hat{a}_{\mathbf{k }s}  $ is an operator that reduce $n_{\mathbf{k }s}$ by one.
Next for the creation operator
\begin{align*}
    \hat{N} \left( \hat{a}^\dagger_{\mathbf{k }s} \ket{n_{\mathbf{k }s}} \right) &=  \hat{a}^\dagger_{\mathbf{k }s} (\hat{N }+1  ) \ket{n_{\mathbf{k }s}}  \\
    \hat{N} \left( \hat{a}^\dagger_{\mathbf{k }s} \ket{n_{\mathbf{k }s}} \right) &=  (n_{\mathbf{k }s}-+  ) \left( \hat{a}^\dagger_{\mathbf{k }s} \ket{n_{\mathbf{k }s}} \right) 
\end{align*}
We use the commutator relation $[AB,C]=A[B,C]-[C,A]$ to obtain the relation
\begin{align*}
  [N,\hat{a}_{\mathbf{k }s} ]&= \hat{a}^\dagger_{\mathbf{k }s} [\hat{a}_{\mathbf{k }s} ,\hat{a}_{\mathbf{k }s} ]+[\hat{a}^\dagger_{\mathbf{k }s} ,\hat{a}_{\mathbf{k }s} ]\hat{a}_{\mathbf{k }s} \\
  \hat{N }\hat{a}_{\mathbf{k }s} -\hat{a}_{\mathbf{k }s} \hat{N}&= \hat{a}_{\mathbf{k }s} \\
  \hat{N }\hat{a}_{\mathbf{k }s} &=  \hat{a}_{\mathbf{k }s} (\hat{N }-1)
\end{align*}
and 
\begin{align*}
  [N,\hat{a}_{\mathbf{k }s}^\dagger ]=\hat{a}_{\mathbf{k }s} [\hat{a}_{\mathbf{k }s} ,\hat{a}_{\mathbf{k }s}^\dagger ]+[\hat{a}^\dagger_{\mathbf{k }s} ,\hat{a}_{\mathbf{k }s}^\dagger ]\hat{a}_{\mathbf{k }s} \\
  \hat{N }\hat{a}_{\mathbf{k }s}^\dagger -\hat{a}_{\mathbf{k }s}^\dagger \hat{N}&= \hat{a}^\dagger_{\mathbf{k }s} \\
  \hat{N }\hat{a}_{\mathbf{k }s}^\dagger = \hat{a}_{\mathbf{k }s}^\dagger (\hat{N }+1)
\end{align*}

Next is to find the coefficients
\begin{equation*}
  \hat{a}_{\mathbf{k }s} \ket{n_{\mathbf{k }s}} = c \ket{n_{\mathbf{k }s}-1}, \qquad
  \hat{a}_{\mathbf{k }s}^\dagger \ket{n_{\mathbf{k }s}} = d \ket{n_{\mathbf{k }s}+1}
\end{equation*}
For the annihilation coefficients
\begin{equation*}
  \braket{n_{\mathbf{k }s} | \hat{N }|n_{\mathbf{k }s}} =
   \left( \hat{a}_{\mathbf{k }s} \ket{n_{\mathbf{k }s}} \right) ^\dagger \hat{a}_{\mathbf{k }s} \ket{n_{\mathbf{k }s}}=|c|^2= n_{\mathbf{k }s}
\end{equation*}
As for creation coefficients
\begin{equation*}
  \braket{n_{\mathbf{k }s} |\hat{a}_{\mathbf{k }s} \hat{a}^\dagger_{\mathbf{k }s} -1 \ket{n_{\mathbf{k }s}}} = \left( \hat{a}^\dagger_{\mathbf{k }s} \ket{n_{\mathbf{k }s}} \right) ^\dagger \hat{a}^\dagger_{\mathbf{k }s} \ket{n_{\mathbf{k }s}}-1=|d|^2-1=n_{\mathbf{k }s}
\end{equation*}

\subsubsection{Derivation: number operator.}
To prove that the eigenvalue of $\hat{N }$ is $n_{\mathbf{k }s}\geq0$ we simply need to consider the norm of $\hat{a}^\dagger_{\mathbf{k }s} \ket{n_{\mathbf{k }}s}$
\begin{equation*}
  ||\hat{a}_{\mathbf{k }s} \ket{n_{\mathbf{k }s}}||=\left( \hat{a}_{\mathbf{k }s} \ket{n_{\mathbf{k }s}} \right) ^\dagger \hat{a}_{\mathbf{k }s} \ket{n_{\mathbf{k }s}}=\braket{n_{\mathbf{k }s} |\hat{N}|n_{\mathbf{k }s} }
\end{equation*}
as something that non-negative.


\subsection{Coherent State}
The coherent state $|\alpha\rangle$, introduced in 1963 by Roy Glauber, is the eigenstate of the annihilation operator $\hat{a}$
\begin{equation*}
  \hat{a}|\alpha\rangle = \alpha|\alpha\rangle
\end{equation*}
with complex eigenvalue $\alpha = |\alpha|e^{i\theta}$.
The complex eigenvalue $\alpha$ can be written as
\begin{equation*}
  \alpha = |\alpha| e^{i\theta}=\sqrt{\bar{N} }e^{i\theta}
\end{equation*}
The value $|\alpha|=\sqrt{\bar{N}}$ is obtained from
\begin{equation*}
  \braket{\hat{N}  } =\braket{\alpha  | \hat{a}^\dagger \hat{a}|\alpha }=|\alpha|^2\implies |\alpha|=\sqrt{\bar{N}}
\end{equation*}
The parameter $\theta$ is the phase of coherent amplitude $\alpha$.
Physically, $\theta$ is the phase of the classical field that the coherent state represents.

$|\alpha\rangle$ does not have a fixed number of photons, i.e. it is not an eigenstate of the number operator $\hat{N}$,
The definition is chosen to make the quantum state most closely resemble classical behavior.
By defining a state as an eigenstate of $\hat{a}$, we select a state where the "field amplitude" itself has a well-defined value.
The reason being that the electromagnetic field can written as a linear combination of the annihilation operators.

Coherent state may be expanded in Fock state as
\begin{equation*}
  |\alpha\rangle = e^{-|\alpha|^2/2} \sum_{n=0}^{\infty} \frac{\alpha^n}{\sqrt{n!}} |n\rangle
\end{equation*}
The expansion demonstrates that the coherent state is not a number operator $\hat{N}$ eigenstate.
In the multimode state, the definition reads
\begin{equation*}
  \hat{a}_{\mathbf{k }s} \ket{\left\{ \alpha_{\mathbf{k }s} \right\} }=\alpha_{\mathbf{k }s} \ket{\left\{ \alpha_{\mathbf{k }s} \right\} }
\end{equation*}
with $  \ket{\left\{ \alpha_{\mathbf{k }s} \right\} } = \bigotimes_{\mathbf{k},s} \ket{\alpha_{\mathbf{k }s}}$, and the Fock state expansion as
\begin{equation*}
  | \{ \alpha_{\mathbf{k }s} \} \rangle
  =
  \exp\left(-\frac{1}{2} \sum_{\mathbf{k},s} |\alpha_{\mathbf{k }s}|^2\right) \sum_{\{ n_{\mathbf{k }s} \}} \left( \prod_{\mathbf{k},s} \frac{\alpha_{\mathbf{k }s}^{n_{\mathbf{k }s}}}{\sqrt{n_{\mathbf{k }s}!}} \right) \ket{\{ n_{\mathbf{k }s} \}}
\end{equation*}

By definition, the electromagnetic field is non-vanishing in this state
\begin{align*}
  \braket{  \hat{\mathbf{E}}  }
      &=   - \sum_{\mathbf{k},s} \sqrt{\frac{2\bar{N}_{\mathbf{k}s}\hbar \omega_k }{ \epsilon_0 V}}  \sin   \left( \mathbf{k} \cdot \mathbf{r} - \omega_k t+\theta_{\mathbf{k}s}  \right) \boldsymbol{\epsilon}_{\mathbf{k}s}\\
  \braket{  \hat{\mathbf{B}}}
 &=   - \sum_{\mathbf{k},s} \sqrt{\frac{2 \bar{N}_{\mathbf{k}s}\hbar \omega_k \mu_0}{V}}  \sin   \left( \mathbf{k} \cdot \mathbf{r} - \omega_k t+\theta_{\mathbf{k}s}  \right) \hat{\mathbf{k}}\times \boldsymbol{\epsilon}_{\mathbf{k}s}
\end{align*}
The energy of electromagnetic field can also be seen to be non-vanishing even at vacuum
\begin{equation*}
    \braket{  \hat{\mathbf{E}}^2 }
    = \braket{\hat{\mathbf{E}} }^2+\sum_{\mathbf{k },s}\frac{\hbar \omega_k }{2\epsilon_0 V}\qquad
  \braket{  \hat{\mathbf{B}}^2 }
    = \braket{\hat{\mathbf{B}} }^2+\sum_{\mathbf{k },s}\frac{\hbar \omega_k}{2 V} \mu_0
\end{equation*}

The coherent state is minimum-uncertainty states, meaning it saturates the Heisenberg bound
\begin{equation*}
  \Delta x \Delta p = \frac{\hbar }{2}
\end{equation*}

\subsubsection{Derivation: Fock state.}
The coherent state $|\alpha\rangle$ can be expanded as
\begin{equation*}
  |\alpha\rangle = \sum_{n=0}^{\infty} c_n |n\rangle
\end{equation*}
Then we have
\begin{equation*}
  \hat{a} |\alpha\rangle = \hat{a} \sum_{n=0}^{\infty} c_n |n\rangle = \sum_{n=0}^{\infty} c_n \hat{a}|n\rangle = \sum_{n=1}^{\infty} c_n \sqrt{n}|n-1\rangle = \sum_{n=0}^{\infty} c_{n+1}\sqrt{n+1} |n\rangle
\end{equation*}
and
\begin{equation*}
  \alpha \ket{\alpha} =  \alpha \sum_{n=0 }^{\infty}c_n \ket{n}= \sum_{n=0 }^{\infty }\alpha c_n \ket{n}
\end{equation*}
Since, by definition $  \hat{a} |\alpha\rangle = \alpha |\alpha\rangle$, it follows that
\begin{align*}
  c_{n+1} \sqrt{n+1} & = \alpha c_n                    \\
  c_{n+1}            & = \frac{\alpha}{\sqrt{n+1}} c_n
\end{align*}
namely
\begin{align*}
  c_1 & =  \frac{\alpha}{\sqrt{1}} c_0                                 \\
  c_2 & =\frac{\alpha }{\sqrt{2}}c_1= \frac{\alpha^2}{\sqrt{2!}} c_0   \\
  c_3 & =  \frac{\alpha }{\sqrt{3 }}c_2=\frac{\alpha^3}{\sqrt{3!}} c_0 \\
      & \vdots
\end{align*}
and in general
\begin{equation*}
  c_n = \frac{\alpha^n}{\sqrt{n!}} c_0
\end{equation*}
Substituting back
\begin{equation*}
  |\alpha\rangle = \sum_{n=0}^{\infty} c_0 \frac{\alpha^n}{\sqrt{n!}} |n\rangle ,
\end{equation*}
where the parameter $c_0$ is fixed by the normalization
\begin{equation*}
  \langle \alpha | \alpha \rangle = |c_0|^2 \sum_{n=0}^{\infty} \frac{\alpha^{2n}}{n!} = |c_0|^2 e^{|\alpha|^2}\implies c_0 = e^{-|\alpha|^2/2}
\end{equation*}

Putting it all together, the single mode state reads
\begin{equation*}
  |\alpha\rangle = e^{-|\alpha|^2/2} \sum_{n=0}^{\infty} \frac{\alpha^n}{\sqrt{n!}} |n\rangle
\end{equation*}

For the multimode coherent state, note that the state is a tensor product of single mode over $(\mathbf{k},s)$
\begin{equation*}
  \ket{\left\{ \alpha_{\mathbf{k }s} \right\} } = \bigotimes_{\mathbf{k},s} \ket{\alpha_{\mathbf{k }s}}=
  \bigotimes_{\mathbf{k },s} e^{-|\alpha|^2/2} \sum_{n=0}^{\infty} \frac{\alpha^n}{\sqrt{n!}} |n\rangle
\end{equation*}
Using the relation
\begin{equation*}
  \bigotimes_i \left(b_i \sum_{n_i } c_{n_i }^{(i )}\ket{n_i}  \right)
  =
  \left( \prod_i b_i \right) \sum_{\left\{ n_i  \right\} }\left( \prod_i  c_{n_i} ^{(i)}\right)  \bigotimes_i \ket{n_i}
\end{equation*}
yields
\begin{align*}
  \ket{\left\{ \alpha_{\mathbf{k }s} \right\} }
   & =
  \left( \prod_{\mathbf{k},s} e^{-|\alpha_{\mathbf{k }s}|^2/2} \right) \sum_{\{n_{\mathbf{k }s}\}} \left( \prod_{\mathbf{k},s} \frac{\alpha_{\mathbf{k }s}^{n_{\mathbf{k }s}}}{\sqrt{n_{\mathbf{k }s}!}} \right) \bigotimes_{\mathbf{k },s}|n_{\mathbf{k }s}\rangle \\
  | \{ \alpha_{\mathbf{k }s} \} \rangle
   & =
  \exp\left(-\frac{1}{2} \sum_{\mathbf{k},s} |\alpha_{\mathbf{k }s}|^2\right) \sum_{\{ n_{\mathbf{k }s} \}} \left( \prod_{\mathbf{k},s} \frac{\alpha_{\mathbf{k }s}^{n_{\mathbf{k }s}}}{\sqrt{n_{\mathbf{k }s}!}} \right) \ket{\{ n_{\mathbf{k }s} \}}
\end{align*}

\subsubsection{Derivation: coherence.}
Form multimode definition
\begin{equation*}
  \sum_{\mathbf{k },s} \hat{a}_{\mathbf{k }s} \ket{\left\{ \alpha_{\mathbf{k }s} \right\} }=\sum_{\mathbf{k },s} \alpha_{\mathbf{k }s} \ket{\left\{ \alpha_{\mathbf{k }s} \right\} }
\end{equation*}
it can be seen that
\begin{align*}
  \sum_{\mathbf{k },s}  \braket{\left\{ \alpha_{\mathbf{k }s } \right\}    | \hat{a}_{\mathbf{k }s} | \left\{ \alpha_{\mathbf{k }s} \right\} }         & =  \sum_{\mathbf{k },s} \alpha_{\mathbf{k }s}   \\
  \sum_{\mathbf{k },s} \braket{\left\{ \alpha_{\mathbf{k }s } \right\}    | \hat{a}^\dagger_{\mathbf{k }s}  | \left\{ \alpha_{\mathbf{k }s} \right\} } & = \sum_{\mathbf{k },s} \alpha_{\mathbf{k }s} ^*
\end{align*}
and
\begin{gather*}
  \sum_{\mathbf{k},s}\sum_{\mathbf{k}',s'} \braket{\hat{a}_{\mathbf{k }s}   \hat{a}_{\mathbf{k }'s'} }
  =\sum_{\mathbf{k},s}\sum_{\mathbf{k}',s'}  \alpha_{\mathbf{k}s}\alpha_{\mathbf{k}'s'} 
  \\
  \sum_{\mathbf{k},s}\sum_{\mathbf{k}',s'} \braket{\hat{a}_{\mathbf{k }s}   \hat{a}^\dagger_{\mathbf{k }'s'} }
  =\sum_{\mathbf{k},s}\sum_{\mathbf{k}',s'} \braket{\hat{a}^\dagger_{\mathbf{k }s} \hat{a}_{\mathbf{k }s} + \delta_{\mathbf{k }\mathbf{k}'}\delta_{ss'}}
  =
    \sum_{\mathbf{k},s} \sum_{\mathbf{k}',s'}\alpha_{\mathbf{k}s}^* \alpha_{\mathbf{k}'s'} 
    +\sum_{\mathbf{k},s}
  \\
  \sum_{\mathbf{k},s}\sum_{\mathbf{k}',s'} \braket{\hat{a}^\dagger_{\mathbf{k }'s'}\hat{a}_{\mathbf{k }s}}
  =
    \sum_{\mathbf{k},s} \sum_{\mathbf{k}',s'}\alpha^{*}_{\mathbf{k}s}  \alpha_{\mathbf{k}'s'} 
  \\
  \sum_{\mathbf{k},s}\sum_{\mathbf{k}',s'} \braket{\hat{a}^\dagger_{\mathbf{k }s}   \hat{a}^\dagger_{\mathbf{k }'s'} }
  = \sum_{\mathbf{k},s} \sum_{\mathbf{k}',s'} \alpha^{*}_{\mathbf{k}s}  \alpha^{*}_{\mathbf{k}'s'} 
\end{gather*}
Also recall Euler identity $  e ^{i\phi}-e ^{-i\phi}=2i \sin \phi$ for good measure.

These relations are enough to derive the following.
First the fact that the electromagnetic
\begin{align*}
  \hat{\mathbf{E}} & =  i \sum_{\mathbf{k},s} \sqrt{\frac{\hbar \omega_k}{2 \epsilon_0 V}} \left[ \hat{a}_{\mathbf{k}s} e^{i\phi_{\mathbf{k}}} - \hat{a}_{\mathbf{k}s}^{\dagger} e^{-i\phi_{\mathbf{k}}} \right] \boldsymbol{\epsilon}_{\mathbf{k}s}                    \\
  \hat{\mathbf{B}} & =  \sum_{\mathbf{k},s} \sqrt{\frac{\hbar \omega_k}{2 V} \mu_0} \left[ \hat{a}_{\mathbf{k}s} e^{i\phi_{\mathbf{k}}} - \hat{a}_{\mathbf{k}s}^{\dagger} e^{-i\phi_{\mathbf{k}}} \right] i \mathbf{\hat{k}} \times \boldsymbol{\epsilon}_{\mathbf{k}s}
\end{align*}
with $\phi_{\mathbf{k }}=\mathbf{k }\cdot \mathbf{r }-\omega_kt$,    expectation value is non-zero in this state.
The expectation value for electric field is derived as follows
\begin{align*}
  \braket{  \hat{\mathbf{E}}  }
   & =  i \sum_{\mathbf{k},s} \sqrt{\frac{\hbar \omega_k}{2 \epsilon_0 V}} \left[ \braket{\hat{a}_{\mathbf{k}s}  } e^{i\phi_{\mathbf{k}}} - \braket{\hat{a}_{\mathbf{k}s}^{\dagger}} e^{-i\phi_{\mathbf{k}}} \right] \boldsymbol{\epsilon}_{\mathbf{k}s} \\
   & =  i \sum_{\mathbf{k},s} \sqrt{\frac{\hbar \omega_k}{2 \epsilon_0 V}} \left[ \alpha_{\mathbf{k }s} e^{i\phi_{\mathbf{k}}} - \alpha_{\mathbf{k }s}^* e^{-i\phi_{\mathbf{k}}} \right] \boldsymbol{\epsilon}_{\mathbf{k}s}                             \\
   & =  i \sum_{\mathbf{k},s} \sqrt{\frac{\bar{N}_{\mathbf{k}s}\hbar \omega_k }{2 \epsilon_0 V}} \left[  e^{i(\mathbf{k} \cdot \mathbf{r} - \omega_k t+ \theta_{\mathbf{k}s})} - e^{-i(\mathbf{k} \cdot \mathbf{r} - \omega_k t+\theta_{\mathbf{k}s}) } \right] \boldsymbol{\epsilon}_{\mathbf{k}s}        \\
  \braket{  \hat{\mathbf{E}}  }
   & =  - \sum_{\mathbf{k},s} \sqrt{\frac{2\bar{N}_{\mathbf{k}s}\hbar \omega_k }{ \epsilon_0 V}}  \sin   \left( \mathbf{k} \cdot \mathbf{r} - \omega_k t+\theta_{\mathbf{k}s}  \right) \boldsymbol{\epsilon}_{\mathbf{k}s}
\end{align*}
Whereas the magnetic field as
\begin{align*}
  \braket{  \hat{\mathbf{B}}  }
   & =   \sum_{\mathbf{k},s} \sqrt{\frac{\hbar \omega_k \mu_0}{2  V}} \left[ \braket{\hat{a}_{\mathbf{k}s}  } e^{i(\mathbf{k} \cdot \mathbf{r} - \omega_k t)} - \braket{\hat{a}_{\mathbf{k}s}^{\dagger}} e^{-i(\mathbf{k} \cdot \mathbf{r} - \omega_k t)} \right] i \hat{\mathbf{k}}\times \boldsymbol{\epsilon}_{\mathbf{k}s}                                          \\
   & =   \sum_{\mathbf{k},s} \sqrt{\frac{\hbar \omega_k\mu_0}{2  V}} \left[ \alpha_{\mathbf{k }s} e^{i(\mathbf{k} \cdot \mathbf{r} - \omega_k t)} - \alpha_{\mathbf{k }s}^* e^{-i(\mathbf{k} \cdot \mathbf{r} - \omega_k t)} \right] i \hat{\mathbf{k}}\times \boldsymbol{\epsilon}_{\mathbf{k}s}                                                                       \\
   & =   \sum_{\mathbf{k},s} \sqrt{\frac{\bar{N}_{\mathbf{k}s}\hbar \omega_k \mu_0}{2 V}} \left[  e^{i(\mathbf{k} \cdot \mathbf{r} - \omega_k t+ \theta_{\mathbf{k}s})} - e^{-i(\mathbf{k} \cdot \mathbf{r} - \omega_k t+\theta_{\mathbf{k}s}) } \right] i \hat{\mathbf{k}}                                                  \times \boldsymbol{\epsilon}_{\mathbf{k}s} \\
  \braket{  \hat{\mathbf{B}}  }
   & =  - \sum_{\mathbf{k},s} \sqrt{\frac{2 \bar{N}_{\mathbf{k}s}\hbar \omega_k \mu_0}{V}}  \sin   \left( \mathbf{k} \cdot \mathbf{r} - \omega_k t+\theta_{\mathbf{k}s}  \right) \hat{\mathbf{k}}\times \boldsymbol{\epsilon}_{\mathbf{k}s}
\end{align*}

Then we move to the square of the field
\begin{multline*}
    \hat{\mathbf{E}}^2
    =  - \sum_{\mathbf{k},s} \sum_{\mathbf{k}',s'} \frac{\hbar \sqrt{\omega_k \omega_{k'}}}{2 \epsilon_0 V}
  \bigg[ \hat{a}_{\mathbf{k }s} \hat{a}_{\mathbf{k }'s'}  e ^{i(\phi_{\mathbf{k }} +\phi_{\mathbf{k }'})} 
  -\hat{a}_{\mathbf{k }s} \hat{a}^\dagger_{\mathbf{k }'s'} e ^{i(\phi_{\mathbf{k }}-\phi_{\mathbf{k }'})} \\ 
   -\hat{a}_{\mathbf{k }s} ^\dagger\hat{a}_{\mathbf{k }'s'} e ^{-i(\phi_{\mathbf{k }}-\phi_{\mathbf{k }'})}  
   +\hat{a}^\dagger_{\mathbf{k }s} \hat{a}^\dagger_{\mathbf{k }'s'} e ^{-i(\phi_{\mathbf{k }}+\phi_{\mathbf{k }'})} \bigg]
  \boldsymbol{\epsilon}_{\mathbf{k}s} \cdot \boldsymbol{\epsilon}_{\mathbf{k}'s'}
\end{multline*}
The expectation value then
\begin{align*}
  \braket{\hat{\mathbf{E}}^2 }
    &=   - \sum_{\mathbf{k},\mathbf{k}', s} \frac{\hbar \sqrt{\omega_k \omega_{k'}}}{2 \epsilon_0 V}
  \bigg[ \braket{\hat{a}_{\mathbf{k }s} \hat{a}_{\mathbf{k }'s'} }  e ^{i(\phi_{\mathbf{k }} +\phi_{\mathbf{k }'})} 
  -\braket{\hat{a}_{\mathbf{k }s} \hat{a}^\dagger_{\mathbf{k }'s'}  } e ^{i(\phi_{\mathbf{k }}-\phi_{\mathbf{k }'})} \\ 
  &\qquad-\braket{\hat{a}^\dagger_{\mathbf{k }s} \hat{a}_{\mathbf{k }'s'} } e ^{-i(\phi_{\mathbf{k }}-\phi_{\mathbf{k }'})}  
   +\braket{\hat{a}^\dagger_{\mathbf{k }s} \hat{a}^\dagger_{\mathbf{k }'s'} } e ^{-i(\phi_{\mathbf{k }}+\phi_{\mathbf{k }'})} \bigg]\\
    &=   - \sum_{\mathbf{k},\mathbf{k}',s} \frac{\hbar \sqrt{\omega_k \omega_{k'}}}{2 \epsilon_0 V}
  \bigg[ \alpha_{\mathbf{k}s} \alpha_{\mathbf{k }' s}e ^{i(\phi_{\mathbf{k }} +\phi_{\mathbf{k }'})} 
 -\alpha_{\mathbf{k }s }\alpha_{\mathbf{k' } s }^* e ^{i(\phi_{\mathbf{k }}-\phi_{\mathbf{k }'})} \\ 
&\qquad-\alpha_{\mathbf{k } s}^*\alpha_{\mathbf{k' } s} e ^{-i(\phi_{\mathbf{k }}-\phi_{\mathbf{k }'})}  
+\alpha_{\mathbf{k } s } ^*\alpha_{\mathbf{k' } s}^* e ^{-i(\phi_{\mathbf{k }}+\phi_{\mathbf{k }'})} -\delta_{\mathbf{k }\mathbf{k}'}\delta_{ss'}e ^{i(\phi_{\mathbf{k }}-\phi_{\mathbf{k }'})} \bigg]\\
    &=   - \sum_{\mathbf{k},\mathbf{k}',s} \frac{\hbar }{2 \epsilon_0 V}\sqrt{\bar{N}_{\mathbf{k }s}\omega_k \bar{N}_{\mathbf{k}' s}\omega_{k'}}
  \bigg[ e ^{i(\phi_{\mathbf{k }} +\phi_{\mathbf{k }'}+\theta_{\mathbf{k }s}+\theta_{\mathbf{k }'s})} \\
&\qquad -e ^{i(\phi_{\mathbf{k }}-\phi_{\mathbf{k }'} +\theta_{\mathbf{k }s}-\theta_{\mathbf{k }' s})}  
 -e ^{-i(\phi_{\mathbf{k }}-\phi_{\mathbf{k }'}-\theta_{\mathbf{k }s}+\theta_{\mathbf{k }' s})}  \\
&\qquad + e ^{-i(\phi_{\mathbf{k }}+\phi_{\mathbf{k }'}-\theta_{\mathbf{k }s}-\theta_{\mathbf{k }'s})} \bigg]+ \sum_{\mathbf{k },s}\frac{\hbar \omega_k }{2\epsilon_0 V}                              \\
    &=   - \sum_{\mathbf{k},\mathbf{k}',s} \frac{2\hbar }{\epsilon_0 V}\sqrt{\bar{N}_{\mathbf{k }s}\omega_k \bar{N}_{\mathbf{k}' s}\omega_{k'}}
    \sin \left( \phi_{\mathbf{k }s}+\theta_{\mathbf{k}s} \right) 
    \sin \left( \phi_{\mathbf{k }'s}+\theta_{\mathbf{k}'s} \right) \\
&\qquad+ \sum_{\mathbf{k },s}\frac{\hbar \omega_k }{2\epsilon_0 V}                              \\
\braket{  \hat{\mathbf{E}}^2 }
   & = \braket{\hat{\mathbf{E}} }^2+\sum_{\mathbf{k },s}\frac{\hbar \omega_k }{2\epsilon_0 V}
\end{align*}
We have used the Euler identity, specifically
\begin{equation*}
    e^{iA} - e^{iB} - e^{-iB} + e^{-iA} = -4 \sin\left(\frac{A + B}{2}\right) \sin\left(\frac{A - B}{2}\right) 
\end{equation*}
which is applied on 
\begin{equation*}
  A = \phi_\mathbf{k} + \phi_{\mathbf{k}'} + \theta_{\mathbf{k}s} + \theta_{\mathbf{k}'s}, B= \phi_\mathbf{k} - \phi_{\mathbf{k}'} + \theta_{\mathbf{k}s} - \theta_{\mathbf{k}'s},
\end{equation*}

Then the magnetic field
\begin{multline*}
    \hat{\mathbf{B}}^2
    =  - \sum_{\mathbf{k},s} \sum_{\mathbf{k}',s'} \frac{\hbar \sqrt{\omega_k \omega_{k'}}}{2 V} \mu_0
  \bigg[ \hat{a}_{\mathbf{k }s} \hat{a}_{\mathbf{k }'s'}  e ^{i(\phi_{\mathbf{k }} +\phi_{\mathbf{k }'})} 
  -\hat{a}_{\mathbf{k }s} \hat{a}^\dagger_{\mathbf{k }'s'} e ^{i(\phi_{\mathbf{k }}-\phi_{\mathbf{k }'})} \\ 
   -\hat{a}_{\mathbf{k }s}^\dagger \hat{a}_{\mathbf{k }'s'}e ^{-i(\phi_{\mathbf{k }}-\phi_{\mathbf{k }'})}  
   +\hat{a}^\dagger_{\mathbf{k }s} \hat{a}^\dagger_{\mathbf{k }'s'} e ^{-i(\phi_{\mathbf{k }}+\phi_{\mathbf{k }'})} \bigg]
  (\mathbf{\hat{k}} \times \boldsymbol{\epsilon}_{\mathbf{k}s}) \cdot (\mathbf{\hat{k}'} \times \boldsymbol{\epsilon}_{\mathbf{k}'s'})
\end{multline*}
which have the expectation value of 
\begin{align*}
  \braket{\hat{\mathbf{B}}^2 }
    &=   - \sum_{\mathbf{k},\mathbf{k}', s} \frac{\hbar \sqrt{\omega_k \omega_{k'}}}{2 V}\mu_0
  \bigg[ \braket{\hat{a}_{\mathbf{k }s} \hat{a}_{\mathbf{k }'s'} }  e ^{i(\phi_{\mathbf{k }} +\phi_{\mathbf{k }'})} 
  -\braket{\hat{a}_{\mathbf{k }s} \hat{a}^\dagger_{\mathbf{k }'s'}  } e ^{i(\phi_{\mathbf{k }}-\phi_{\mathbf{k }'})} \\ 
  &\qquad-\braket{\hat{a}^\dagger_{\mathbf{k }s} \hat{a}_{\mathbf{k }'s'} } e ^{-i(\phi_{\mathbf{k }}-\phi_{\mathbf{k }'})}  
   +\braket{\hat{a}^\dagger_{\mathbf{k }s} \hat{a}^\dagger_{\mathbf{k }'s'} } e ^{-i(\phi_{\mathbf{k }}+\phi_{\mathbf{k }'})} \bigg]\\
    &=   - \sum_{\mathbf{k},\mathbf{k}',s} \frac{\hbar \sqrt{\omega_k \omega_{k'}}}{2 V}\mu_0  \bigg[ \alpha_{\mathbf{k}s} \alpha_{\mathbf{k }' s}e ^{i(\phi_{\mathbf{k }} +\phi_{\mathbf{k }'})} 
 -\alpha_{\mathbf{k }s }\alpha_{\mathbf{k' } s }^* e ^{i(\phi_{\mathbf{k }}-\phi_{\mathbf{k }'})} \\ 
&\qquad-\alpha_{\mathbf{k } s}^*\alpha_{\mathbf{k' } s} e ^{-i(\phi_{\mathbf{k }}-\phi_{\mathbf{k }'})}  
+\alpha_{\mathbf{k } s } ^*\alpha_{\mathbf{k' } s}^* e ^{-i(\phi_{\mathbf{k }}+\phi_{\mathbf{k }'})} -\delta_{\mathbf{k }\mathbf{k}'}\delta_{ss'}e ^{i(\phi_{\mathbf{k }}-\phi_{\mathbf{k }'})} \bigg]\\
    &=   - \sum_{\mathbf{k},\mathbf{k}',s} \frac{\hbar }{2 V}\sqrt{\bar{N}_{\mathbf{k }s}\omega_k \bar{N}_{\mathbf{k}' s}\omega_{k'}}\mu_0
  \bigg[ e ^{i(\phi_{\mathbf{k }} +\phi_{\mathbf{k }'}+\theta_{\mathbf{k }s}+\theta_{\mathbf{k }'s})} \\
&\qquad -e ^{i(\phi_{\mathbf{k }}-\phi_{\mathbf{k }'} +\theta_{\mathbf{k }s}-\theta_{\mathbf{k }' s})}  
 -e ^{-i(\phi_{\mathbf{k }}-\phi_{\mathbf{k }'}-\theta_{\mathbf{k }s}+\theta_{\mathbf{k }' s})}  \\
&\qquad + e ^{-i(\phi_{\mathbf{k }}+\phi_{\mathbf{k }'}-\theta_{\mathbf{k }s}-\theta_{\mathbf{k }'s})} \bigg]+ \sum_{\mathbf{k },s}\frac{\hbar \omega_k }{2V}\mu_0                              \\
    &=   - \sum_{\mathbf{k},\mathbf{k}',s} \frac{2\hbar }{ V}\sqrt{\bar{N}_{\mathbf{k }s}\omega_k \bar{N}_{\mathbf{k}' s}\omega_{k'}}\mu_0
    \sin \left( \phi_{\mathbf{k }s}+\theta_{\mathbf{k}s} \right) 
    \sin \left( \phi_{\mathbf{k }'s}+\theta_{\mathbf{k}'s} \right) \\
&\qquad+ \sum_{\mathbf{k },s}\frac{\hbar \omega_k }{2V}\mu_0                              \\
\braket{  \hat{\mathbf{B}}^2 }
   & = \braket{\hat{\mathbf{B}} }^2+\sum_{\mathbf{k },s}\frac{\hbar \omega_k }{2V}\mu_0
\end{align*}

We can insert this into classical Hamiltonian
\begin{align*}
  H & =   \frac{1}{ 2}\int d^3 \mathbf{r} \sum_{\mathbf{k},s} \left( \epsilon_0\mathbf{E}^2 + \frac{1}{\mu_0} \mathbf{B}^2 \right)        \\
    & = \frac{1 }{2} \int d^3 \mathbf{r } \sum_{\mathbf{k},s} \left[ \epsilon_0\braket{\mathbf{E} }^2+ \frac{1 }{\mu_0}\braket{\mathbf{B} }^2 +\frac{\hbar \omega_k }{V}\right] \\
    & = \frac{1 }{2} \int d^3 \mathbf{r } \sum_{\mathbf{k},\mathbf{k}', s} \bigg[ \frac{4\hbar }{V}\sqrt{\bar{N}_{\mathbf{k }s}\omega_k \bar{N}_{\mathbf{k}' s}\omega_{k'}}
    \sin \left( \phi_{\mathbf{k }s}+\theta_{\mathbf{k}s} \right) 
    \sin \left( \phi_{\mathbf{k }'s}+\theta_{\mathbf{k}'s} \right)\\ 
    &\qquad+\frac{\hbar \omega_k }{V}\bigg] \\
  H & = \sum_{\mathbf{k },s}\left[ \bar{N }_{\mathbf{k }s}\hbar \omega_k+\frac{\hbar \omega_k }{2} \right]=  \sum_{\mathbf{k},s} \hbar \omega_k \left( \bar{N}_{\mathbf{k }s}+\frac{1 }{2} \right)
\end{align*}
to recover the same result of quantum Hamiltonian
\begin{align*}
  \braket{\left\{ \alpha_{\mathbf{k }s} \right\} |H|\left\{ \alpha_{\mathbf{k }s} \right\} }
   & =
  \left\langle \left\{ \alpha_{\mathbf{k }s} \right\} \Bigg| \sum_{\mathbf{k},s} \hbar \omega_k \left( \hat{N }_{\mathbf{k} s }+ \frac{1 }{2} \right)\Bigg|\left\{ \alpha_{\mathbf{k }s} \right\}  \right\rangle \\
  \braket{H}
   & =  \sum_{\mathbf{k},s} \hbar \omega_k \left( \bar{N}_{\mathbf{k }s} +\frac{1 }{2} \right)
\end{align*}

\subsubsection{Derivation: minimum uncertainty.}
The variance of $x$ is 
\begin{align*}
  (\Delta x)^2 &=  \braket{x^2 } - \braket{x }^2 \\
  &= \frac{\hbar}{2m\omega} \braket{a^2 + a^{\dagger 2} + aa^\dagger + a^\dagger a  } 
  - \left( \sqrt{\frac{\hbar}{2m\omega}} \langle a + a^{\dagger} \rangle \right) ^2\\
  &= \frac{\hbar}{2m\omega} \braket{a^2 + a^{\dagger 2} + aa^\dagger + a^\dagger a }- \frac{\hbar}{2m\omega} (\alpha + \alpha^{*})^2 \\
  &=\frac{\hbar}{2 m \omega} \left( \alpha^{2} + \alpha^{*2} + 2 |\alpha|^{2} + 1 \right)  - \frac{\hbar}{2 m \omega} \left( \alpha^2 + \alpha^{*2} + 2 |\alpha|^2 \right)\\
  \Delta x
  &= \sqrt{\frac{\hbar  }{2m\omega}}
\end{align*}
In other hand, the variance of $p$ is
\begin{align*}
  (\Delta p)^2 &=  \braket{p^2 } - \braket{p }^2 \\
  &= - \frac{m\hbar\omega}{2} \braket{a^2 + a^{\dagger 2} -aa^\dagger - a^\dagger a  } 
  - \left( -i\sqrt{\frac{m\hbar\omega}{2}} \langle a - a^{\dagger} \rangle \right) ^2\\
  &= - \frac{m\hbar\omega}{2} \braket{a^2 + a^{\dagger 2} -2 a^\dagger a-1 } + \frac{m\hbar\omega}{2} (\alpha - \alpha^{*})^2 \\
  &= - \frac{ m \hbar  \omega}{2 } \left( \alpha^{2} + \alpha^{*2} - 2 |\alpha|^{2} -1 \right)  + \frac{m \hbar  \omega}{2 } \left( \alpha^2 + \alpha^{*2} - 2 |\alpha|^2 \right)\\
  \Delta p
  &= \sqrt{\frac{m  \hbar \omega }{2}}
\end{align*}
multiplying both gives the minimum uncertainty.
\end{document}